\section{Introduction}
\label{sec:intro}
Large language models (LLMs) have evolved beyond text to handle a wide variety of multimodal tasks~\citep{cui2025emu3,ming_flash,wu2025qwen,gu2025ui,liu2026lumina}.
%
Understanding and generation represent the two primary categories of multimodal tasks. 
%
By modeling both visual understanding and generation as token sequence prediction, LLMs have achieved highly competitive results in both areas. 
%
Traditionally, these tasks are handled by separate specialized models, such as Qwen-VL~\citep{bai2025qwen2,team2025qwen3VL} or InternVL~\citep{Internvl,internvl2,internvl35} for understanding, and Flux~\citep{flux} or Z-Image~\citep{cai2025z} for generation. 
%
However, a unified model that handles both within a single framework offers several key benefits: it promotes mutual enhancement between understanding and generation, improves deployment efficiency, and unlocks advanced capabilities like interleaved generation and reasoning, ultimately bringing us closer to artificial general intelligence (AGI).


Current unified multimodal models predominantly build upon autoregressive (AR) architectures.
%
Janus~\citep{Janus} and Lumina-mGPT~\citep{liu2026lumina} tokenizes images into discrete sequences and unifies both modalities under next-token prediction, while OmniGen2~\citep{wu2025omnigen2}, Hunyuan Image 3.0~\citep{cao2025hunyuanimage}, and BAGEL~\citep{bagel} adopt a hybrid paradigm combining text autoregression with image diffusion.
%
While these AR-based approaches have shown promise, masked diffusion models~\citep{lou2023discrete,sahoo2024simple,xin2025resurrect} offer an alternative paradigm with inherent advantages in parallel decoding and bidirectional context modeling.
A unified masked diffusion framework further simplifies training through a single objective, avoiding the delicate balance between AR and diffusion losses.
%
However, existing unified masked diffusion models, such as MMaDA~\citep{Mmada} and Lumina-DiMOO~\citep{dimoo}, still lag behind state-of-the-art AR-based unified architectures in both task coverage and benchmark performance.
%
This gap fundamentally stems from their architecture and modeling designs: 1) their reconstructive VQ tokenizers lack semantic information, causing poor understanding performance; 2) excessive image compression by VQ tokenizers compromises generation quality; 3) their fully bidirectional modeling has been shown to be unreliable for text.
%
Furthermore, they commonly assume fixed output lengths for understanding tasks, limiting their applicability in open-ended scenarios.

To overcome these limitations, we propose LLaDA2.0-Uni, a unified dLLM-based Mixture-of-Experts (MoE) model for seamless multimodal understanding and generation.
%
At its core, LLaDA2.0-Uni utilizes LLaDA2.0~\citep{LLaDA2} (a 16B dLLM MoE architecture) as its backbone.
%
A key architectural innovation is the introduction of the SigLIP-VQ tokenizer, which converts continuous visual inputs into fully discrete semantic tokens.
%
Unlike previous reconstruction-based tokenizers that struggle with multimodal understanding, this purely semantic representation preserves crucial details and effectively supports complex visual reasoning.
%
Consequently, this design maintains the unified discrete modeling format, allowing both text and images to be optimized under a shared block-level masked diffusion objective.
%
For image generation, LLaDA2.0-Uni employs a dedicated Diffusion Decoder to process the discrete tokens generated by the dLLM backbone.
%
Optimized through distillation, this decoder synthesizes high-fidelity images in just 8 inference steps, achieving an excellent balance between speed and quality.

LLaDA2.0-Uni achieves top-tier performance across both understanding and generation benchmarks, as shown in Figure~\ref{fig:main}.
%
In multimodal understanding, LLaDA2.0-Uni demonstrates competitive visual question answering and document reasoning capabilities compared with specialized VLMs such as Qwen2.5-VL~\citep{bai2025qwen2}
%
Regarding image generation, LLaDA2.0-Uni produces high-quality images and enables highly flexible image editing.
%
Beyond these general tasks, the unified discrete representation empowers LLaDA2.0-Uni to support interleaved generation and reasoning.
%
This flexibility establishes LLaDA2.0-Uni as a powerful and efficient paradigm for the next generation of unified foundation models.


The key contributions of LLaDA2.0-Uni can be summarized as follows:

\begin{itemize}[leftmargin=16pt]
    \item \textbf{Novel Unified Architecture.} LLaDA2.0-Uni integrates a fully semantic tokenizer, a 16B MoE dLLM backbone, and a diffusion decoder. This architecture unifies text and image modeling through a shared block-wise mask prediction objective.
    
    \item \textbf{Interleaved Generation and Reasoning.} Beyond its strong performance in both understanding and generation, LLaDA2.0-Uni inherently supports interleaved generation and reasoning, marking a significant step toward exploring how generation and understanding can reinforce each other.
    
    \item \textbf{Efficient Inference.} Building on the advantages of parallel decoding, LLaDA2.0-Uni further accelerates inference by optimizing the decoding process in the dLLM backbone and applying few-step distillation to the decoder, achieving an effective balance between speed and performance.

    \item \textbf{Strong Benchmark Performance.} LLaDA2.0-Uni achieves strong performance across visual understanding, generation, and editing benchmarks, performing on par with state-of-the-art unified models.
\end{itemize}

